\documentclass{article}
\usepackage{hyperref}
\usepackage{Style}

\nocite{*} % Comentar si quiero citar
%\addbibresource{bibliografia.bib} % Quitar el comentado si quiero usar bibliografia

\begin{document}

\begin{minipage}{2.5cm}
    \includegraphics[width=2cm]{imagen_puc.jpg}
\end{minipage}
\begin{minipage}{12cm}
    {\sc Pontificia Universidad Católica de Chile\\
    Facultad de Matemáticas\\
    Departamento de Matemática\\
    Profesor: Manuel Cabezas -- Ayudante: Benjamín Mateluna}
\end{minipage}
\vspace{2ex}

{\centerline{\bf Introducción a la Geometría - MAT1304}
\centerline{\bf Ayudantía 19 - Repaso I5}}
\centerline{\bf 28 de mayo de 2026}

\vspace{1cm}
\noindent\textbf{Problema 1.} Resuelva la ecuación $z^{4}=-1+\sqrt{3}i$ en $\C$.

\vspace{5mm}
\noindent\textbf{Problema 2.} Muestre que $p(x),q(x)\in\R[x]$ tienen un factor en común en $\R[x]$
si y solo si tienen una raíz compleja en común.

\vspace{5mm}
\noindent\textbf{Problema 3.} Muestre que $x^{4}+x^{3}-x^{2}+x+1$ es irreducible en $\Q[x]$.

\vspace{5mm}
\noindent\textbf{Problema 4.} Factorice completamente
\begin{enumerate}
    \item $p(x)=x^{4}-1$ en $\R[x]$ y $\C[x]$.
    \item $p(x)=x^{4}+1$ en $\R[x]$ y $\C[x]$.
\end{enumerate}

\vspace{1cm}
\noindent\textbf{\underline{Ejercicios Propuestos:}}

\vspace{5mm}
\noindent\textbf{Extra 1.} Calcule todos los valores que puede tomar la expresión 
$(3-3i)\sqrt{1+i}$.

\vspace{5mm}
\noindent\textbf{Extra 2.} Muestre que $p(x)$ es reducible si y solo si $p(x-a)$ es reducible 
para algún $a\in k$, con $k=\Q,\R,\C$.

%\printbibliography % Quitar el comentado si quiero usar bibliografia

\end{document}